% !TeX root = ../papers/nc-p3.tex
\section{Ore型構成の内部完全性と楕円同型不変量}\label{sec:ore-complete}
\subsection{次数一のOre導分の完全な分類}\label{subsec:ore-derivations}
固定された$B_{t,h}$の上でのOre型構成がどの程度一般的かを、
具体的な導分空間の分類によって明らかにする。ここで$R=k[t,h,(1+th)^{-1}]$、
$\sigma_B|_{(B_{t,h})_1}=-1$とする。
$\operatorname{Der}_1(B)$は内部次数1の通常の導分、
$\operatorname{Der}_{\sigma_B,1}(B)$は内部次数1の$\sigma_B$-導分の$R$-加群とする。

\begin{theorem}[固定Clifford部分のOre導分の全分類]\label{thm:ore-derivations}
$R$上、次の自由加群としての等式が成立する。
\begin{align}
 \operatorname{Der}_1(B_{t,h})
 &=R\operatorname{ad}_W\oplus R\operatorname{ad}_Y\oplus R\operatorname{ad}_Z
   \oplus Rd_a\oplus Rd_b\oplus Rd_c,\label{eq:ordinary-derivations}\\
 \operatorname{Der}_{\sigma_B,1}(B_{t,h})
 &=\operatorname{Hom}_R\bigl((B_{t,h})_1,
       RW^2\oplus RY^2\oplus RZ^2\bigr)\simeq R^9.
       \label{eq:sign-derivations}
\end{align}
ここで$d_a,d_b,d_c$は式\eqref{eq:elliptic-derivation}の$a,b,c$の係数である。
従って、$Xj+jX=p_j+(-1)^i d_j$で表される二周期Ore型拡大で、
$p$を次数1の$\sigma_B$-導分、$d$を通常の次数1の導分とするものは、
固定した$B_{t,h}$上では$9+6=15$個の線型係数を持つ。
$X_i\mapsto X_i+\ell+(-1)^i m$（$\ell,m\in(B_{t,h})_1$）による
六つの独立な冗長性を除くと、残る係数は$6+3=9$個である。
この完全性は\emph{指定したClifford部分代数とOre形式を保持する構成}に限る。
一般の四次元AS正則$\mathbb Z$-代数全体の分類を意味しない。
\end{theorem}
\begin{proof}
$(B_{t,h})_2$の基底を$(W^2,WY,WZ,Y^2,YZ,Z^2)$とする。
補題\ref{lem:ternary-clifford}の中心平方上の表示によって、
三次成分には
\[
 W(W^2,Y^2,Z^2),\quad Y(W^2,Y^2,Z^2),\quad
 Z(W^2,Y^2,Z^2),\quad WYZ
\]
の10個の基底がある。生成元3個の各像を六つの基底で書くと18未知数となる。
$YW+WY-tZ^2,ZW+WZ-hY^2,ZY+YZ-W^2$への
Leibniz則または$\sigma_B$-Leibniz則の適用は、
これらの三次基底上に30行18列の行列$M_{\rm ord},M_\sigma$を与える。

通常導分の行列には、$t,h$に依存しない行列式$-2048$を持つ
12次正方小行列がある。よって$\operatorname{rank}M_{\rm ord}\ge12$。
一方、三つの内導分と式\eqref{eq:elliptic-derivation}の三つの導分は核に属し、
それらの18行6列の係数行列には、行列式$4$の6次正方小行列がある。
従って全ファイバーで核の次元は少なくとも6であり、上の12次小行列とあわせて
厳密に6となる。二つの定数小行列式が単元であるため、$R$上でも
六つの導分が核の自由基底となる。

$\sigma_B$-導分については、三つの生成元の像を任意の中心二次式とすれば
三関係への適用が零となるので、$R^9$が核に含まれる。
一方$M_\sigma$には、二つの9次正方小行列があり、その行列式は
\[
 4(th-2)(th+4)^2,\qquad 3ht^3(th-8)^2
\]
である。これらは同時に零にならないため、全ファイバーで階数は9以上である。
中心二次式を取る九つの標準導分は直和因子であり、
二つの小行列が覆う開集合上で$M_\sigma$の核と一致する。
ゆえに$R$上でも式\eqref{eq:sign-derivations}を得る。

最後に$\ell,m\in(B_{t,h})_1$として生成元を変えると、
関係の右辺への変化は
\[
 p_j\longmapsto p_j+\ell j+j\ell,\qquad
 d_j\longmapsto d_j+mj-jm
\]
である。前者は三次元の中心二次式値の$\sigma_B$-内導分、
後者は三次元の通常の内導分であり、互いに独立である。
実際$\ell=\alpha W+\beta Y+\gamma Z$とすれば、
$p_W$の$W^2$係数、$p_Z$の$W^2,Z^2$係数への寄与は
$2\alpha,\beta,2\gamma$なので、前者は全ファイバーで自由に作用する。
後者の独立性は上の六導分の自由基底性から従う。
全行列、小行列の行・列番号、厳密な多項式恒等式は
\path{computations/verify_ore_derivations.py}で再現できる。
\end{proof}

\begin{corollary}[九変数正規形のOre型での十分性]\label{cor:ore-slice}
$B_{t,h}$と符号自己同型$\sigma_B$を固定した上の二周期Ore型構成を考える。
$X_i\mapsto X_i+\ell+(-1)^im$という変更を許せば、
$d$は一意に$d_a,d_b,d_c$の線型結合へ移せる。
また$p$は$p_W$の$W^2$係数と$p_Z$の$W^2,Z^2$係数を零にした六係数の形へ移せる。

さらに$B_{t,h}$の三関係の形を保つ対角変換
$W\mapsto\lambda W,Y\mapsto\mu Y,Z\mapsto\nu Z,X\mapsto\rho X$
（$\lambda^2=\mu\nu$）を許し、残る$p_Y$の$W^2,Z^2$係数が非零な開集合に限ると、
両係数を1にできる。従ってこの\emph{印付きOre型構成}の稠密開集合は、
有限個の対角正規化の選択を除いて式\eqref{eq:nine-family}の九変数族で表される。
従って同じClifford部分のOre形式を固定したまま、さらに自由な
次数一の導分係数を付け加えても、一般の同型類の次元は増えない。
\end{corollary}
\begin{proof}
定理\ref{thm:ore-derivations}の六つの基底変更によって所要の切片を取る。
前者三個は通常の内導分を除き、後者三個は中心二次式の三係数を除く。
$\lambda^2=\mu\nu$を満たす対角変換は、全体スカラーを除いて二次元のトーラスをなす。
$p_Y$の$W^2,Z^2$係数を$A_0,B_0$とすると、変換後は
$\rho A_0/\nu,\rho\mu B_0/\nu^2$である。
従って$A_0B_0\ne0$なら$\rho/\nu=A_0^{-1}$、
$\mu/\nu=A_0/B_0$と取って両者を1にできる。
$\lambda^2=\mu\nu$に残る符号の有限選択だけは同型類の次元に影響しない。
残る六つの中心係数は$(u,v,r,s)$および規格化された二係数、
通常導分の外部三係数は$(a,b,c)$へ整理され、
式\eqref{eq:nine-family}と一致する。$t,h$も動かした全パラメータ数は
$2+15-3-3-2=9$である。
これは印付きを忘れた同型の完全分類ではない。
九変数族の\emph{裸の}indexed同型類の像の次元9は、前節の独立な階数証明による。
\end{proof}

\subsection{内在的な楕円$j$-不変量による同型の必要条件}
\begin{proposition}[九変数族の同型に対する明示的障害]\label{prop:elliptic-j-obstruction}
九変数族のうち、$E_{t,h}$が滑らかで、唯一の非被約左点が$e_X$、
その点からの射影で残りの左点が$E_{t,h}$上の相異なる12点となる開集合を考える。
この開集合の二つのファイバーが\emph{添字保持}$\ZZ$-代数として同型なら、
$x=th$、$x'=t'h'$に対して
\begin{equation}\label{eq:elliptic-j}
 J(x)=J(x'),\qquad
 J(x)=\frac{(x-2)^3(x^3-6x^2+228x-8)^3}
                 {x(x-8)^3(x+1)^6}
\end{equation}
が必要である。$J:\mathbb P^1_x\to\mathbb P^1_j$は次数12の非定数有理写像である。
従って同型類を一つ固定したとき、その同型類に属し得る一般の$x=th$の値は
高々12個である。ただし$j$の一致は同型の\emph{十分条件ではない}。
\end{proposition}
\begin{proof}
添字保持同型は左点スキームの射影同型を誘導する。
その唯一の非被約点$e_X$は内在的であるから射影同型で保たれる。
従ってそこからの射影が作る平面上の12点集合も射影同型で対応する。
それらを通る平面三次曲線は$E_{t,h}$ただ一つである。
実際、滑らかな既約三次曲線と異なる三次曲線は、共通成分がなければ
B\'ezoutの定理により9点を超えて相異なる点で交わらない。
よって二つのファイバーが同型なら$E_{t,h}$と$E_{t',h'}$が同型となり、
それらの$j$-不変量は一致する。

標数零の代数閉体上で$w,y,z$をそれぞれ適切に三乗根で規格化すると、
式\eqref{eq:elliptic-cubic}はHesse標準形
$W^3+Y^3+Z^3-3\kappa WYZ=0$になり、
\[
 \kappa^3=-\frac{(2-th)^3}{27th}.
\]
Hesse三次曲線の古典的公式
$j=27\kappa^3(\kappa^3+8)^3/(\kappa^3-1)^3$
（\cite[\S7.3, (42)]{KLV}）へ代入して式\eqref{eq:elliptic-j}を得る。
分子の次数12、分母の次数10で、両者に共通因子はない。
従って有理写像の次数は12であり、各幾何学的ファイバーに含まれる
$x$の値は高々12個である。
多項式恒等式と次数の検証は\path{computations/verify_ore_derivations.py}にある。
\end{proof}
